\textbf{COLUNI 2019 -- Questão 03}\par\medskip \begin{center}\begin{verse}
No descomeço era o verbo.\\
Só depois é que veio o delírio do verbo.\\
O delírio do verbo estava no começo, lá onde a\\
criança diz: \textit{Eu escuto a cor dos passarinhos.}\\
A criança não sabe que o verbo escutar não funciona\\
para cor, mas para som.\\
Então se a\\
criança muda a função de um verbo, ele\\
delira. E pois.\\
Em poesia que é voz de poeta, que é a voz de fazer\\
nascimentos ―\\
O verbo tem que pegar delírio.
\end{verse}
(BARROS, Manoel de. \textbf{Meu quintal é maior do que o mundo}. Rio de Janeiro: Alfaguara, 2015, p. 83.)\end{center}

Ao utilizar a palavra “descomeço” em lugar de começo, o eu lírico:\begin{enumerate}
\item afirma que a criação da palavra poética se deu por meio da palavra de Deus.
\item recupera o princípio da simultaneidade da criação do objeto a partir da palavra.
\item situa a origem e existência da palavra poética em uma dimensão atemporal.
\item integra a concretização da palavra de Deus e da palavra poética.\end{enumerate}
